\chapter{Help systems}
\label{ch:help}

\begin{aside}
A help screen is a TUI's documentation. Make it
context-aware, searchable, and accessible without leaving
the app.
\end{aside}

\section{Triggering help}

Standard keybind: F1 or \code{?}. Some apps use
\code{:help} (vim).

\section{Modal vs split}

A help overlay (modal) covers the screen briefly. A help
panel (split) stays open beside the work. Pick based on
the app's complexity.

\section{Context-aware}

\code{?} in normal mode shows normal-mode help.
\code{?} in a dialog shows dialog-specific help. The
current mode or focus drives the help content.

\section{Content sources}

\begin{itemize}[leftmargin=*]
  \item Keymap auto-generated (Chapter~\ref{ch:keymaps}).
  \item Markdown files shipped with the app.
  \item Tooltip text from individual widgets.
\end{itemize}

\section{Search}

A search box in the help screen lets the user find
``how do I rename''. The result list shows matching
sections.

\section{Glossary}

A glossary of terms with cross-references. Hyperlink
between entries (OSC-8 internal links).

\section{Tutorials}

An interactive tutorial guides the user through tasks.
\maya{}'s \code{tutorial} widget shows steps and tracks
progress.

\section{Recap}

\begin{recap}
\begin{itemize}[leftmargin=*]
  \item F1 or \code{?} to open.
  \item Context-aware content.
  \item Sources: keymap, markdown, tooltips.
  \item Search, glossary, tutorials.
\end{itemize}
\end{recap}

\section*{Workshop}

\begin{enumerate}[leftmargin=*]
  \item Build a help overlay that lists keybindings.
  \item Make it context-aware (different content per
    mode).
  \item Add a search box that filters help content.
\end{enumerate}
